% Source: https://artofproblemsolving.com/community/c3h2800473_2020_osn_mathematics_city_level_questions_16_grade_89
% Stable identity: aops-topic-2800473-q03
$f(x+1,x-2,5)=(x+1)(x-2)+5(x-2)+5(x+1)=x^2+9x-7$. The equation $x^2+9x-7=3x^2-12$ becomes $2x^2-9x-5=0$, i.e. $(2x+1)(x-5)=0$. The larger value of $2x_1-3x_2$ is with $x_1=5$, $x_2=-\dfrac12$: $10+\dfrac32=\boxed{\textbf{(d) }\dfrac{23}2}$.
